\section{Plan de pruebas}

En este apartado, se describen las pruebas desarrolladas para cada uno de los submódulos individuales, así como para el PCS completo que integra los tres submódulos.
Estas fueron implementadas para asegurar el correcto funcionamiento de las máquinas de estados separadas, así como en conjunto.
Se probó una variedad de casos para cada uno, donde se contempló el funcionamiento normal y algunas variaciones.

\subsection{Módulo \texttt{TRANSMIT}}

En esta subsección se describen las pruebas funcionales aplicadas al módulo de transmisión completo, utilizando el probador en \texttt{src/transmit/tester.v}.
Cada prueba define una secuencia de octetos de datos (\texttt{txd}) y una activación de la señal \texttt{tx\_en}, con el objetivo de verificar el comportamiento del transmisor respecto a:
\begin{itemize}
  \item La generación correcta de los \textit{ordered sets}: \texttt{/S/}, \texttt{/D/}, \texttt{/T/}, \texttt{/R/} e \texttt{/I/}.
  \item La presencia o ausencia de \textit{carrier extend} (\texttt{/R R/} después de \texttt{/T/}).
  \item El manejo del \textit{running disparity} y la selección del IDLE correcto (\texttt{/I1/} o \texttt{/I2/}).
\end{itemize}

A continuación se detalla cada una de las pruebas definidas.

\subsubsection*{Prueba 0: IDLEs después de \texttt{reset}}

\begin{itemize}
  \item \textbf{Descripción:} Tras el \texttt{reset}, se espera un período de inactividad donde \texttt{tx\_en = 0} y el transmisor permanece en estado de \texttt{IDLE}.
  \item \textbf{Resultado esperado:}
    \begin{itemize}
      \item El transmisor debe generar únicamente \textit{ordered sets} de tipo \texttt{/I/} (secuencias \texttt{/K28.5/} seguidas de \texttt{/D16.2/}.
      \item No deben aparecer símbolos de inicio de trama (\texttt{/S/}), datos (\texttt{/D/}), terminación (\texttt{/T/}) ni extendidos (\texttt{/R/}) mientras \texttt{tx\_en = 0}.
    \end{itemize}
\end{itemize}

\subsubsection*{Prueba 1: Carrier extend y IDLE con disparidad correcta}

\begin{itemize}
  \item \textbf{Descripción:} Se habilita la transmisión (\texttt{tx\_en = 1}) y se envía la siguiente secuencia de datos:
    \[
      \texttt{D11\_3, D23\_1, D7\_4, D12\_5}
    \]
    Posteriormente, \texttt{tx\_en} se desactiva (\texttt{tx\_en = 0}), lo que fuerza el envío del \textit{ordered set} de terminación \texttt{/T/}. Según el estándar y la lógica implementada en \texttt{transmit\_ordered\_set}, esta configuración debe producir \emph{carrier extend}, pues el primer /R/ está en ciclo par.
  \item \textbf{Resultado esperado:}
    \begin{itemize}
      \item La secuencia de \textit{ordered sets} debe ser:
      \[
        \texttt{/S/ D11\_3 D23\_1 D7\_4 D12\_5 /T/ /R/ /R/}
      \]
      es decir, se observa \texttt{/T/} seguido de dos \texttt{/R/}, indicando \textit{carrier extend}.
      \item Tras el \textit{carrier extend}, el transmisor debe retornar al envío de IDLEs \texttt{/I/} con un \emph{running disparity} negativo.
    \end{itemize}
\end{itemize}

\subsubsection*{Prueba 2: Sin carrier extend e IDLE con disparidad errónea}

\begin{itemize}
  \item \textbf{Descripción:} Se habilita la transmisión y se envía la secuencia:
    \[
      \texttt{D11\_3, D23\_1, D7\_4, D12\_5, D8\_6}
    \]
    Luego se desactiva \texttt{tx\_en}. Esta combinación de datos está seleccionada para que el \emph{running disparity} al final de la trama sea \textbf{positivo}, pero sin requerir carrier extend.
  \item \textbf{Resultado esperado:}
    \begin{itemize}
      \item La secuencia de \textit{ordered sets} debe ser:
      \[
        \texttt{/S/ D11\_3 D23\_1 D7\_4 D12\_5 D8\_6 /T/ /R/}
      \]
      \item Como el \textit{running disparity} final de la trama es positivo, el siguiente IDLE debe ser \texttt{/I1/}, compuesto por:
        \[
          \texttt{/K28.5/ seguido de /D5.6/}
        \]
        con el objetivo de corregir la disparidad y devolverla al valor negativo deseado.
    \end{itemize}
\end{itemize}

\subsubsection*{Prueba 3: Sin carrier extend e IDLE con disparidad correcta}

\begin{itemize}
  \item \textbf{Descripción:} Se habilita la transmisión y se envía la secuencia:
    \[
      \texttt{D11\_3, D23\_1, D1\_0, D31\_1}
    \]
    para luego desactivar \texttt{tx\_en}. En este caso, la combinación de datos se escogió para que el \textit{running disparity} al final de la trama sea \textbf{negativo}.
  \item \textbf{Resultado esperado:}
    \begin{itemize}
      \item La secuencia de \textit{ordered sets} debe ser:
      \[
        \texttt{/S/ D11\_3 D23\_1 D1\_0 D31\_1 /T/ /R/}
      \]
      nuevamente sin \textit{carrier extend}.
      \item Como el \textit{running disparity} final es negativo, el IDLE que sigue debe ser \texttt{/I2/}, compuesto por:
        \[
          \texttt{/K28.5/ seguido de /D16.2/}
        \]
        lo cual mantiene la disparidad en el valor negativo especificado por el estándar.
    \end{itemize}
\end{itemize}


\subsection{Módulo \texttt{SYNCHRONIZATION}}

\subsubsection*{Prueba 1: Sincronización inicial sin errores}

\begin{itemize}
  \item \textbf{Descripción:}  
    Tras el reinicio, se envían cuatro secuencias IDLE del tipo \texttt{/I2/}, cada una compuesta por:
    \[
      \texttt{/K28.5/} \rightarrow \texttt{/D16.2/}
    \]
    con \texttt{indicate = 1} en cada ciclo. Esta secuencia contiene el carácter de \texttt{COMMA}, por lo que la PCS debe usarla para adquirir sincronización.
  \item \textbf{Resultado esperado:}
    \begin{itemize}
      \item El bloque debe entrar en estado sincronizado después de recibir suficientes \texttt{COMMA}, lo cual equivale a  \texttt{code\_sync\_status=1}.
    \end{itemize}
\end{itemize}

\subsubsection*{Prueba 2: Un dato inválido}

\begin{itemize}
  \item \textbf{Descripción:}  
    Tras estar sincronizado, se ingresa un único code-group inválido:
    \[
      \texttt{11\_1111\_1111}
    \]
    seguido inmediatamente de una secuencia válida:
    \[
      \texttt{D5\_6, D16\_2, D0\_0, D2\_0, D5\_6, D16\_2}
    \]
  \item \textbf{Resultado esperado:}
    \begin{itemize}
      \item La lógica debe detectar el inválido y pasar al estado de \texttt{SYNC\_ACQUIRED\_2}.
      \item Al recibir nuevamente 4 code-groups válidos, la PCS debe \textbf{mantener la sincronización}, ya que un único error no supera el umbral de pérdida y vuelve a \texttt{SYNC\_ACQUIRED\_1}.
    \end{itemize}
\end{itemize}

\subsubsection*{Prueba 3: Dos inválidos consecutivos}

\begin{itemize}
  \item \textbf{Descripción:}  
    Se envían dos code-groups inválidos:
    \[
      \texttt{11\_1111\_1111,\ 00\_0000\_0001}
    \]
    seguidos de una secuencia de code-groups válidos del tipo:
    \[
      \texttt{D5\_6, D16\_2, D0\_0, D2\_0, \dots}
    \]
  \item \textbf{Resultado esperado:}
    \begin{itemize}
      \item Transición al estado \texttt{SYNC\_ACQUIRED\_3}.
      \item Los code-groups válidos posteriores deben permitir que la PCS vuelva al estado estable \texttt{SYNC\_ACQUIRED\_1}.
    \end{itemize}
\end{itemize}

\subsubsection*{Prueba 4: Tres inválidos consecutivos}

\begin{itemize}
  \item \textbf{Descripción:}  
    Se envían tres code-groups inválidos seguidos:
    \[
      \texttt{11\_1111\_1111,\ 00\_0000\_0001,\ 11\_1111\_1111}
    \]
    y luego una secuencia extensa de code-groups válidos.
  \item \textbf{Resultado esperado:}
    \begin{itemize}
      \item El sistema debe acercarse al umbral de pérdida de sincronización (4 inválidos).
      \item La secuencia posterior de 12 code-groups válidos debe permitir recuperar la sincronía.
    \end{itemize}
\end{itemize}

\subsubsection*{Prueba 5: Inválidos intercalados durante resincronización}

\begin{itemize}
  \item \textbf{Descripción:}  
    Se envía primero un inválido:
    \[
      \texttt{11\_1111\_1111}
    \]
    luego comienzan a recibirse code-groups válidos, pero durante ese proceso se inyecta otro inválido intercalado:
    \[
      \texttt{00\_0000\_0001}
    \]
  \item \textbf{Resultado esperado:}
    \begin{itemize}
      \item La lógica no debe perder sincronía prematuramente.
      \item El sistema debe continuar volver a \texttt{SYNC\_ACQUIRED\_1} después de recibir 8 code-groups válidos seguidos.
    \end{itemize}
\end{itemize}

\subsubsection*{Prueba 6: Pérdida completa de sincronización}

\begin{itemize}
  \item \textbf{Descripción:}  
    Se colocan cuatro code-groups inválidos, alternando patrones incorrectos:
    \[
      \texttt{11\_1111\_1111,\ 00\_0000\_0001,\ 11\_1111\_1111,\ 00\_0000\_0001,\dots}
    \]
  \item \textbf{Resultado esperado:}
    \begin{itemize}
      \item El contador de errores debe superar el umbral definido por la arquitectura y pasar al estado \texttt{LOSS\_OF\_SYNC}.
      \item La PCS debe declarar pérdida de sincronización y esperar nuevos COMMA para adquirirla de nuevo.
    \end{itemize}
\end{itemize}

\subsection{Módulo \texttt{RECEIVE}}

Para el caso del \texttt{RECEIVE}, se realizaron tres pruebas para corroborar el correcto funcionamiento de este módulo.

\subsubsection*{Prueba 1: Envío de datos con receptor sincronizado}

\begin{itemize}
  \item \textbf{Descripción:}  
    El receptor entra en estado sincronizado mediante la activación de la señal \texttt{sync\_status}.  
    Una vez sincronizado, recibe el ordered set \texttt{/S/} y procede con la decodificación de los datos, enviándolos posteriormente a la interfaz GMII.  
    Los datos recibidos durante esta prueba corresponden a los mostrados en el Cuadro \ref{prueba1_receive}.

    \begin{table}[H]
        \centering
        \begin{tabular}{ccc}
        \hline
            Dato de 10 bits recibido & Running Disparity & Valor octeto\\ \hline
             110100\,1100 & Negativo & 6B\\
             111010\,1001 & Negativo & 37\\ 
             000111\,0010 & Positivo & 87\\ 
             111001\,0110 & Negativo & C8\\ \hline
        \end{tabular}
        \caption{Datos de 10 bits recibidos en la prueba 1.}
        \label{prueba1_receive}
    \end{table}

  \item \textbf{Resultado esperado:}
    \begin{itemize}
      \item El receptor debe decodificar correctamente cada code-group de acuerdo con su disparidad.
      \item Los valores octeto deben enviarse al GMII respetando el orden y manteniendo la coherencia de disparidad según el estándar.
      \item Termina en /T/R/;
    \end{itemize}
\end{itemize}

% -------------------------------------------------------

\subsubsection*{Prueba 2: Envío de datos con receptor no sincronizado}

\begin{itemize}
  \item \textbf{Descripción:}  
    En esta prueba, la señal \texttt{sync\_status} no se encuentra activa, por lo que el receptor no está sincronizado.  
    Los code-groups recibidos nuevamente corresponden al Cuadro \ref{prueba1_receive}, pero al no existir sincronización válida, el receptor no debe entregar datos a la interfaz GMII.

  \item \textbf{Resultado esperado:}
    \begin{itemize}
      \item Aunque se reciba la misma secuencia de code-groups que en la prueba anterior, no debe generarse salida hacia GMII.
      \item El receptor debe permanecer en estado no-sincronizado y descartar los datos recibidos, cumpliendo con el estándar.
    \end{itemize}
\end{itemize}

% -------------------------------------------------------

\subsubsection*{Prueba 3: Terminación incorrecta /T/R/R/}

\begin{itemize}
  \item \textbf{Descripción:}  
    En condiciones normales, la terminación de la trama debe ser \texttt{/T/R/I/}.  
    En esta prueba se fuerza un escenario incorrecto donde la terminación recibida es \texttt{/T/R/R/}.  
    Nuevamente, los datos previos recibidos corresponden al Cuadro \ref{prueba1_receive}.

  \item \textbf{Resultado esperado:}
    \begin{itemize}
      \item El receptor debe detectar que la secuencia \texttt{/T/R/R/} no cumple con el final válido de trama según 1000BASE-X.
      \item Debe irse al estado TRR extend en donde ahí rxd tomará el valor de 0x0F y se activa la señal de rx er.
      \item Seguido a esto concluirá normalmente y se devuelve a esperar un secuencia de envío de datos.
    \end{itemize}
\end{itemize}

\subsection{PCS completa}

Ahora bien, en esta sección se muestra el plan de pruebas desarrollado para los tres módulos conectados entre sí en modo de \textit{loopback}, como lo solicita el enunciado.


\subsubsection*{Prueba 0: IDLEs y adquisición de sincronización}

\begin{itemize}
  \item \textbf{Descripción:}  
    Con \texttt{tx\_en = 0} y mediante la tarea \texttt{ciclos(9)}, el transmisor envía únicamente \texttt{/I/} (IDLEs).  
    Esta fase permite que el sincronizador del PCS detecte los \textit{COMMA} y adquiera sincronización.
  \item \textbf{Resultado esperado:}
    \begin{itemize}
      \item \texttt{code\_sync\_status = 1}.
      \item No deben generarse errores en \texttt{rx\_er} ni datos válidos en \texttt{rx\_dv}.
    \end{itemize}
\end{itemize}

\subsubsection*{Prueba 1: Carrier extend e IDLE con disparidad correcta}

\begin{itemize}
  \item \textbf{Descripción:}  
    Con \texttt{tx\_en = 1} se envía la trama:
    \[
      \texttt{D11\_3,\ D23\_1,\ D7\_4,\ D12\_5}
    \]
    y luego se desactiva \texttt{tx\_en}.  
    La longitud del paquete requiere insertar \texttt{/R/} como \emph{carrier extend} para que el siguiente /K28.5/ quede en ciclo par; además, la disparidad final permite que el IDLE resultante sea correcto.
  \item \textbf{Resultado esperado:}
    \begin{itemize}
      \item \texttt{rx\_dv} se activa únicamente durante la transmisión de datos.
      \item La secuencia de fin de trama debe ser \texttt{/T/ /R/ /R/} (carrier extend).
      \item Se activa \textbf{\texttt{rx\_er}} en el último \texttt{/R/}, como corresponde al final del carrier extend.
      \item La disparidad del IDLE posterior es la correcta según el \emph{running disparity}.
    \end{itemize}
\end{itemize}

\subsubsection*{Prueba 2: Sin carrier extend e IDLE con disparidad incorrecta}

\begin{itemize}
  \item \textbf{Descripción:}  
    Se transmite la trama:
    \[
      \texttt{D11\_3,\ D23\_1,\ D7\_4,\ D12\_5,\ D8\_6}
    \]
    cuyas cinco codificaciones producen un \emph{running disparity} final que genera un IDLE \textbf{incorrecto}.
  \item \textbf{Resultado esperado:}
    \begin{itemize}
      \item \textbf{No hay carrier extend}: la secuencia termina con \texttt{/T/ /R/}.
      \item La PCS debe escoger el IDLE equivocado según la disparidad obtenida, lo cual permite validar el caso de disparidad errónea (/D5.6/).
      \item \texttt{rx\_dv} se activa durante los datos y se apaga al cerrar la trama.
    \end{itemize}
\end{itemize}


\subsubsection*{Prueba 3: Sin carrier extend e IDLE con disparidad correcta}

\begin{itemize}
  \item \textbf{Descripción:}
    Se envía una trama corta:
    \[
      \texttt{D11\_3,\ D23\_1,\ D31\_1}
    \]
    que finaliza con una disparidad coherente con el IDLE siguiente sin necesidad de carrier extend.
  \item \textbf{Resultado esperado:}
    \begin{itemize}
      \item \textbf{No hay carrier extend}: la terminación es solamente \texttt{/T/ /R/}.
      \item La disparidad del IDLE posterior es correcta (/D16.2/ después del /K28.5/.
      \item \texttt{rx\_dv} se activa únicamente por los tres octetos válidos.
    \end{itemize}
\end{itemize}


\subsubsection*{Prueba 4: Envío del subconjunto completo de datos definidos}

\begin{itemize}
  \item \textbf{Descripción:}  
    Se transmite un conjunto representativo amplio de datos:
    \[
      \texttt{D5\_6,\ D16\_2,\ D0\_0,\ D1\_0,\dots,\ D4\_6}
    \]
    cubriendo todos los casos contemplados en la codificación 8b/10b (subconjunto del original).
  \item \textbf{Resultado esperado:}
    \begin{itemize}
      \item La decodificación es correcta y \texttt{rxd} reproduce todos los valores en el orden enviado.
      \item \texttt{rx\_dv} permanece alto durante toda la transmisión.
    \end{itemize}
\end{itemize}


\subsubsection*{Prueba 5: Pérdida completa de sincronización}

\begin{itemize}
  \item \textbf{Descripción:}  
    Se transmiten cinco octetos inválidos:
    \[
      \texttt{0x81,\ 0x81,\ 0x81,\ 0x81,\ 0x81}
    \]
    Como aclaración, no es que los datos sean inválidos, pues el estándar presenta la codificación para los 256 datos de 8 bits. Sin embargo, para simular un dato inválido, se utilizaron datos que no están dentro del conjunto de code-groups definidos.
    En una implementación real, este caso de simulación no funcionaría, pero permite que se detecten code-groups en el sincronizador que no existen.
  \item \textbf{Resultado esperado:}
    \begin{itemize}
      \item Se pierde la sincronización: \texttt{code\_sync\_status = 0}.
      \item Los datos no deben entregarse a través de \texttt{rxd}.
      \item \texttt{rx\_dv} permanece en cero porque no hay datos válidos.
      \item Se pueden activar errores en \texttt{rx\_er}, dependiendo del diseño.
    \end{itemize}
\end{itemize}


